\documentclass[a4paper,12pt,twoside]{memoir}

% Castellano
\usepackage[spanish,es-tabla]{babel}
\selectlanguage{spanish}
\usepackage[utf8]{inputenc}
\usepackage[T1]{fontenc}
\usepackage{lmodern} % scalable font
\usepackage{microtype}
\usepackage{placeins}

\RequirePackage{booktabs}
\RequirePackage[table]{xcolor}
\RequirePackage{xtab}
\RequirePackage{multirow}

% Links
\PassOptionsToPackage{hyphens}{url}\usepackage[colorlinks]{hyperref}
\hypersetup{
	allcolors = {red}
}

% Ecuaciones
\usepackage{amsmath}

% Rutas de fichero / paquete
\newcommand{\ruta}[1]{{\sffamily #1}}

% Párrafos
\nonzeroparskip

% Huérfanas y viudas
\widowpenalty100000
\clubpenalty100000

% Evitar solapes en el header
\nouppercaseheads

% Imagenes
\usepackage{graphicx}
\newcommand{\imagen}[2]{
	\begin{figure}[!h]
		\centering
		\includegraphics[width=0.9\textwidth]{#1}
		\caption{#2}\label{fig:#1}
	\end{figure}
	\FloatBarrier
}

\newcommand{\imagenflotante}[2]{
	\begin{figure}%[!h]
		\centering
		\includegraphics[width=0.9\textwidth]{#1}
		\caption{#2}\label{fig:#1}
	\end{figure}
}



% El comando \figura nos permite insertar figuras comodamente, y utilizando
% siempre el mismo formato. Los parametros son:
% 1 -> Porcentaje del ancho de página que ocupará la figura (de 0 a 1)
% 2 --> Fichero de la imagen
% 3 --> Texto a pie de imagen
% 4 --> Etiqueta (label) para referencias
% 5 --> Opciones que queramos pasarle al \includegraphics
% 6 --> Opciones de posicionamiento a pasarle a \begin{figure}
\newcommand{\figuraConPosicion}[6]{%
  \setlength{\anchoFloat}{#1\textwidth}%
  \addtolength{\anchoFloat}{-4\fboxsep}%
  \setlength{\anchoFigura}{\anchoFloat}%
  \begin{figure}[#6]
    \begin{center}%
      \Ovalbox{%
        \begin{minipage}{\anchoFloat}%
          \begin{center}%
            \includegraphics[width=\anchoFigura,#5]{#2}%
            \caption{#3}%
            \label{#4}%
          \end{center}%
        \end{minipage}
      }%
    \end{center}%
  \end{figure}%
}

%
% Comando para incluir imágenes en formato apaisado (sin marco).
\newcommand{\figuraApaisadaSinMarco}[5]{%
  \begin{figure}%
    \begin{center}%
    \includegraphics[angle=90,height=#1\textheight,#5]{#2}%
    \caption{#3}%
    \label{#4}%
    \end{center}%
  \end{figure}%
}
% Para las tablas
\newcommand{\otoprule}{\midrule [\heavyrulewidth]}
%
% Nuevo comando para tablas pequeñas (menos de una página).
\newcommand{\tablaSmall}[5]{%
 \begin{table}
  \begin{center}
   \rowcolors {2}{gray!35}{}
   \begin{tabular}{#2}
    \toprule
    #4
    \otoprule
    #5
    \bottomrule
   \end{tabular}
   \caption{#1}
   \label{tabla:#3}
  \end{center}
 \end{table}
}

%
%Para el float H de tablaSmallSinColores
\usepackage{float}

%
% Nuevo comando para tablas pequeñas (menos de una página).
\newcommand{\tablaSmallSinColores}[5]{%
 \begin{table}[H]
  \begin{center}
   \begin{tabular}{#2}
    \toprule
    #4
    \otoprule
    #5
    \bottomrule
   \end{tabular}
   \caption{#1}
   \label{tabla:#3}
  \end{center}
 \end{table}
}

\newcommand{\tablaApaisadaSmall}[5]{%
\begin{landscape}
  \begin{table}
   \begin{center}
    \rowcolors {2}{gray!35}{}
    \begin{tabular}{#2}
     \toprule
     #4
     \otoprule
     #5
     \bottomrule
    \end{tabular}
    \caption{#1}
    \label{tabla:#3}
   \end{center}
  \end{table}
\end{landscape}
}

%
% Nuevo comando para tablas grandes con cabecera y filas alternas coloreadas en gris.
\newcommand{\tabla}[6]{%
  \begin{center}
    \tablefirsthead{
      \toprule
      #5
      \otoprule
    }
    \tablehead{
      \multicolumn{#3}{l}{\small\sl continúa desde la página anterior}\\
      \toprule
      #5
      \otoprule
    }
    \tabletail{
      \hline
      \multicolumn{#3}{r}{\small\sl continúa en la página siguiente}\\
    }
    \tablelasttail{
      \hline
    }
    \bottomcaption{#1}
    \rowcolors {2}{gray!35}{}
    \begin{xtabular}{#2}
      #6
      \bottomrule
    \end{xtabular}
    \label{tabla:#4}
  \end{center}
}

%
% Nuevo comando para tablas grandes con cabecera.
\newcommand{\tablaSinColores}[6]{%
  \begin{center}
    \tablefirsthead{
      \toprule
      #5
      \otoprule
    }
    \tablehead{
      \multicolumn{#3}{l}{\small\sl continúa desde la página anterior}\\
      \toprule
      #5
      \otoprule
    }
    \tabletail{
      \hline
      \multicolumn{#3}{r}{\small\sl continúa en la página siguiente}\\
    }
    \tablelasttail{
      \hline
    }
    \bottomcaption{#1}
    \begin{xtabular}{#2}
      #6
      \bottomrule
    \end{xtabular}
    \label{tabla:#4}
  \end{center}
}

%
% Nuevo comando para tablas grandes sin cabecera.
\newcommand{\tablaSinCabecera}[5]{%
  \begin{center}
    \tablefirsthead{
      \toprule
    }
    \tablehead{
      \multicolumn{#3}{l}{\small\sl continúa desde la página anterior}\\
      \hline
    }
    \tabletail{
      \hline
      \multicolumn{#3}{r}{\small\sl continúa en la página siguiente}\\
    }
    \tablelasttail{
      \hline
    }
    \bottomcaption{#1}
  \begin{xtabular}{#2}
    #5
   \bottomrule
  \end{xtabular}
  \label{tabla:#4}
  \end{center}
}



\definecolor{cgoLight}{HTML}{EEEEEE}
\definecolor{cgoExtralight}{HTML}{FFFFFF}

%
% Nuevo comando para tablas grandes sin cabecera.
\newcommand{\tablaSinCabeceraConBandas}[5]{%
  \begin{center}
    \tablefirsthead{
      \toprule
    }
    \tablehead{
      \multicolumn{#3}{l}{\small\sl continúa desde la página anterior}\\
      \hline
    }
    \tabletail{
      \hline
      \multicolumn{#3}{r}{\small\sl continúa en la página siguiente}\\
    }
    \tablelasttail{
      \hline
    }
    \bottomcaption{#1}
    \rowcolors[]{1}{cgoExtralight}{cgoLight}

  \begin{xtabular}{#2}
    #5
   \bottomrule
  \end{xtabular}
  \label{tabla:#4}
  \end{center}
}




\graphicspath{ {./img/} }

% Capítulos
\chapterstyle{bianchi}
\newcommand{\capitulo}[2]{
	\setcounter{chapter}{#1}
	\setcounter{section}{0}
	\setcounter{figure}{0}
	\setcounter{table}{0}
	\chapter*{#2}
	\addcontentsline{toc}{chapter}{#2}
	\markboth{#2}{#2}
}

% Apéndices
\renewcommand{\appendixname}{Apéndice}
\renewcommand*\cftappendixname{\appendixname}

\newcommand{\apendice}[1]{
	%\renewcommand{\thechapter}{A}
	\chapter{#1}
}

\renewcommand*\cftappendixname{\appendixname\ }

% Formato de portada
\makeatletter
\usepackage{xcolor}
\newcommand{\tutor}[1]{\def\@tutor{#1}}
\newcommand{\course}[1]{\def\@course{#1}}
\definecolor{cpardoBox}{HTML}{E6E6FF}
\def\maketitle{
  \null
  \thispagestyle{empty}
  % Cabecera ----------------
\noindent\includegraphics[width=\textwidth]{cabecera}\vspace{1cm}%
  \vfill
  % Título proyecto y escudo informática ----------------
  \colorbox{cpardoBox}{%
    \begin{minipage}{.8\textwidth}
      \vspace{.5cm}\Large
      \begin{center}
      \textbf{TFG del Grado en Ingeniería Informática}\vspace{.6cm}\\
      \textbf{\LARGE\@title{}}
      \end{center}
      \vspace{.2cm}
    \end{minipage}

  }%
  \hfill\begin{minipage}{.20\textwidth}
    \includegraphics[width=\textwidth]{escudoInfor}
  \end{minipage}
  \vfill
  % Datos de alumno, curso y tutores ------------------
  \begin{center}%
  {%
    \noindent\LARGE
    Presentado por \@author{}\\ 
    en Universidad de Burgos --- \@date{}\\
    Tutor: \@tutor{}\\
  }%
  \end{center}%
  \null
  \cleardoublepage
  }
\makeatother


% Datos de portada
\title{Draw E-R App \\Documentación Técnica}
\author{Rubén Maté Iturriaga}
\tutor{Jesús Manuel Maudes Raedo}
\date{\today}

\begin{document}

\maketitle



\cleardoublepage



%%%%%%%%%%%%%%%%%%%%%%%%%%%%%%%%%%%%%%%%%%%%%%%%%%%%%%%%%%%%%%%%%%%%%%%%%%%%%%%%%%%%%%%%



\frontmatter


\clearpage

% Indices
\tableofcontents

\clearpage

\listoffigures

\clearpage

\listoftables

\clearpage

\mainmatter

\appendix

\apendice{Plan de Proyecto Software}

\section{Introducción}
En este apartado se realiza un pequeño resumen del desarrollo del proyecto y de los diferentes sprints que lo han compuesto, así como un estudio de viabilidad en el que se presenta tanto la viabilidad económica como legal del proyecto.

\section{Planificación temporal}
Para el desarrollo de este trabajo de fin de grado se ha usado una metodología de desarrollo ágil. Los sprints han sido de aproximadamente unas 3-4 semanas, lo cuál ha permitido poder compaginarlo con mi trabajo.
Se han administrado con las Milestones de la plataforma GitHub en combinación con las labels que determinan el tipo de tarea que es cada issue.

\imagen{sprints}{División en sprints del desarrollo}
\imagen{labels}{Labels a asignar a los diferentes issues y pull requests}

\subsection{Sprint 0 - 13/03/2024 - 8/04/2024 }
El primer sprint ha consistido en plantear las bases sobre las que se iba a cimentar el proyecto. Esto es:
\begin{itemize}
    \item Decidir el framework tecnológico a usar. El debate estaba entre React \cite{react} o Javascript \cite{wiki:javascript}, se optó finalmente por este primero puesto que no se le ve un gran sentido a desarrollar puramente con Javascript a estas alturas del desarrollo web.
    \item Decidir la librería para dibujar/modelar: Se optó finalmente por mxGraph\cite{mxgraph}/maxGraph\cite{maxgraph}.
\end{itemize}

\imagen{sprint-0}{Issues del Sprint 0}

\begin{table}[ht!]
    \centering
    \resizebox{14cm}{!} {
    \begin{tabular}{|l|c|}
    \hline
    \rowcolor[rgb]{0.9, 0.9, 0.9}
    \textbf{Issues} & \textbf{Label} \\ \hline
    \textbf{Decision: Canvas Drawing Library Selection} & \cellcolor[rgb]{0.83, 0.77, 0.97}research \\ \hline 
    \textbf{Consideration: Framework Selection - React vs Vanilla Javascript} & \cellcolor[rgb]{0.83, 0.77, 0.97}research \\ \hline 
    \end{tabular}}
    \caption{Issues completadas Sprint 0}
\end{table}

\subsection{Sprint 1 - 9/04/2024 - 24/04/2024 }
El segundo sprint ha consistido en:
\begin{itemize}
    \item Decidir qué formato de datos va a usarse para guardar y poder tratar un diagrama E-R. Este issue no se completará totalmente hasta una fase casi final del proyecto.
    \item Añadir la plantilla de la memoria LaTeX al repositorio.
    \item Implementar la funcionalidad inicial de poder añadir Entidades y cambiar su nombre.
\end{itemize}

\imagen{sprint-1}{Issues del Sprint 1}

\begin{table}[ht!]
    \centering
    \resizebox{14cm}{!} {
    \begin{tabular}{|l|c|}
    \hline
    \rowcolor[rgb]{0.9, 0.9, 0.9}
    \textbf{Issues} & \textbf{Label} \\ \hline
    \textbf{Decision: Choose Data Format for Saving E-R Diagrams} & \cellcolor[rgb]{0.83, 0.77, 0.97}research \\ \hline 
    \textbf{Documentation: Add LaTeX Final Degree Project Report} & \cellcolor[rgb]{0.0, 0.46, 0.79}\textcolor{white}{documentation} \\ \hline 
    \textbf{Feature: Add Interlinked Entities and Text Labels to Canvas} & \cellcolor[rgb]{0.64, 0.93, 0.94}enhancement \\ \hline 
    \end{tabular}}
    \caption{Issues completadas Sprint 1}
\end{table}


\subsection{Sprint 2 - 25/04/2024 - 29/05/2024 }
El tercer sprint ha consistido en:
\begin{itemize}
    \item Establecer un framework de testing, tanto para pruebas unitarias como end to end.
    \item Establecer un pipeline de integración y despliegue continuos haciendo uso de las Actions de GitHub \cite{github-actions}.
    \item Añadir atributos a las entidades.
    \item Mostrar/ocultar atributos de las entidades con un menú contextual.
    \item Añadir una función de utilidad que compruebe si un determinado diagrama es válido. Esta función no se completará en su totalidad hasta un punto posterior.
\end{itemize}

\imagen{sprint-2}{Issues del Sprint 2}

\begin{table}[ht!]
    \centering
    \resizebox{14cm}{!} {
    \begin{tabular}{|l|c|}
    \hline
    \rowcolor[rgb]{0.9, 0.9, 0.9}
    \textbf{Issues} & \textbf{Label} \\ \hline
    \textbf{Tooling: Add CI/CD pipeline} & \cellcolor[rgb]{0.0, 0.42, 0.46}\textcolor{white}{testing} \\ \hline 
    \textbf{Tooling: Add a testing framework} & \cellcolor[rgb]{0.0, 0.42, 0.46}\textcolor{white}{testing} \\ \hline 
    \textbf{Feature: Add attributes to the entities} & \cellcolor[rgb]{0.64, 0.93, 0.94}enhancement \\ \hline 
    \textbf{Feature: Show/hide entities attributes with a contextual menu} & \cellcolor[rgb]{0.64, 0.93, 0.94}enhancement \\ \hline 
    \textbf{Feature: Add a utility function to check if a certain graph is valid} & \cellcolor[rgb]{0.64, 0.93, 0.94}enhancement \\ \hline 
    \end{tabular}}
    \caption{Issues completadas Sprint 2}
\end{table}


\subsection{Sprint 3 - 29/05/2024 - 15/06/2024 }
El cuarto sprint ha consistido en:
\begin{itemize}
    \item Cambiar la iconografía alrededor de los atributos de las entidades para que el atributo clave tuviera su label subrayada.
    \item Cambiar el modo en que se añaden los atributos. Ya no se escoge si se quiere añadir clave, sino que el primer atributo añadido es clave y el resto no. Con un menú contextual se puede convertir no clave en clave.
    \item Añadir relaciones (simples) entre entidades.
\end{itemize}

\imagen{sprint-3}{Issues del Sprint 3}

\begin{table}[ht!]
    \centering
    \resizebox{14cm}{!} {
    \begin{tabular}{|l|c|}
    \hline
    \rowcolor[rgb]{0.9, 0.9, 0.9}
    \textbf{Issues} & \textbf{Label} \\ \hline
    \textbf{Feature: Underline key attribute} & \cellcolor[rgb]{0.64, 0.93, 0.94}enhancement \\ \hline 
    \textbf{Feature: Primary (key) and secondary attributes management via contextual menu} & \cellcolor[rgb]{0.64, 0.93, 0.94}enhancement \\ \hline 
    \textbf{Feature: Add (simple) relations between entities} & \cellcolor[rgb]{0.64, 0.93, 0.94}enhancement \\ \hline 
    \end{tabular}}
    \caption{Issues completadas Sprint 3}
\end{table}

\subsection{Sprint 4 - 15/06/2024 - 05/07/2024 }
El cuarto y último sprint ha consistido en:
\begin{itemize}
    \item Implementar la funcionalidad para que las relaciones 'N:M' puedan tener atributos.
    \item Implementar la funcionalidad para prevenir que atributos, entidades y relaciones se inicien con nombres repetidos.
    \item Implementar la funcionalidad para generar el script SQL correspondiente al paso a tablas del diagrama modelado.
    \item Implementar la funcionalidad para añadir botones de Borrar en los diversos elementos.
    \item Implementar la funcionalidad para poder reconfigurar relaciones.
    \item Implementar la funcionalidad para que los atributos se muevan al mover la entidad y evitar que los nuevos atributos se inicialicen todos apilados.
    \item Implementar la funcionalidad para guardar el diagrama en el almacenamiento local del navegador y poder importar/exportar los diagramas en formato JSON.
\end{itemize}

\imagen{sprint-4-1}{Issues del Sprint 4 (Parte 1)}
\imagen{sprint-4-2}{Issues del Sprint 4 (Parte 2)}

\begin{table}[ht!]
    \centering
    \resizebox{14cm}{!} {
    \begin{tabular}{|l|c|}
    \hline
    \rowcolor[rgb]{0.9, 0.9, 0.9}
    \textbf{Issues} & \textbf{Label} \\ \hline
    \textbf{Feature: 'N:M' relations can have attributes} & \cellcolor[rgb]{0.64, 0.93, 0.94}enhancement \\ \hline 
    \textbf{Feature: Prevent Duplicate Names in Attributes, Entities and Relations} & \cellcolor[rgb]{0.64, 0.93, 0.94}enhancement \\ \hline 
    \textbf{Feature: Generate SQL script for the tables of the current diagram} & \cellcolor[rgb]{0.64, 0.93, 0.94}enhancement \\ \hline 
    \textbf{Feature: Add delete buttons for the different elements} & \cellcolor[rgb]{0.64, 0.93, 0.94}enhancement \\ \hline 
    \textbf{Feature: Reconfigure relations} & \cellcolor[rgb]{0.64, 0.93, 0.94}enhancement \\ \hline
    \textbf{Feature: Better attributes drawing management} & \cellcolor[rgb]{0.64, 0.93, 0.94}enhancement \\ \hline
    \textbf{Feature: Save the diagram on localhost, add export/import diagram} & \cellcolor[rgb]{0.64, 0.93, 0.94}enhancement \\ \hline
    \textbf{Documentation} & \cellcolor[rgb]{0.0, 0.46, 0.79}\textcolor{white}{documentation} \\ \hline
    \end{tabular}}
    \caption{Issues completadas Sprint 4}
\end{table}

\newpage
\section{Estudio de viabilidad}
En esta sección se va a realizar un estudio de la viabilidad económica y legal del proyecto.

\subsection{Viabilidad económica}

\subsubsection{Costes}
En primer lugar, se detallan los costes asociados al proyecto.

\begin{itemize}
\item \textbf{Costes de personal}

En este apartado se detalla el gasto asociado al personal involucrado en el proyecto. 
La estimación que se ha hecho de horas de trabajo es de 400 horas repartidas a lo largo de 4 meses, con un cálculo de 25 horas semanales. El salario del alumno se estima en 20 EUR/hora, lo que se traduce en un salario bruto mensual de:

$$ 25\frac{horas}{semana} \times 20\frac{\text{EUR}}{hora} \times 4\frac{semanas}{mes} = 2000 \text{ EUR} \hspace{0.5em}al\hspace{0.5em}mes $$

Este cálculo corresponde al salario bruto del empleado.

Para determinar el salario neto que recibirá el empleado, se deben incluir los impuestos que la empresa debe pagar por él. Estos impuestos se han consultado y obtenido de la página oficial de la seguridad social: \url{http://www.seg-social.es/wps/portal/wss/internet/Trabajadores/CotizacionRecaudacionTrabajadores/36537}.

\begin{table}[ht!]
    \centering
    \begin{tabular}{|l|c|}
    \hline
    \textbf{Tipo de impuesto} & \textbf{Porcentaje} \\ \hline
    Contingencias & 23.6\% \\ \hline
    Desempleo & 5.5\% \\ \hline
    FOGASA & 0.20\% \\ \hline
    Formación profesional & 0.60\% \\ \hline
    \end{tabular}
    \caption{Porcentajes de impuestos aplicables}
    \label{tab:impuestos}
\end{table}

\imagen{seg-social}{Régimen general de la Seguridad Social}

Teniendo en cuenta estos impuestos, el coste total del empleado es:

$$\frac{2000\frac{\text{ EUR}}{mes}}{1 - (0.236 + 0.055 + 0.002 + 0.006)} = 2.855,77 \text{ EUR} \hspace{0.5em}al\hspace{0.5em}mes$$

Además, se cuenta con un profesor contratado como apoyo, con un salario de 35 EUR/hora debido a su experiencia. El profesor trabaja tres horas a la semana:

$$ 3\frac{horas}{semana} \times 35\frac{\text{EUR}}{hora} \times 4\frac{semanas}{mes} = 420 \text{ EUR} \hspace{0.5em}al\hspace{0.5em}mes $$

Sumando el profesor, se obtiene un total de 420 EUR brutos, a los cuales también se deben sumar los impuestos:

$$\frac{420\frac{\text{ EUR}}{mes}}{1 - (0.236 + 0.055 + 0.002 + 0.006)} = 599,18 \text{ EUR} \hspace{0.5em}al\hspace{0.5em}mes \hspace{0.5em}para\hspace{0.5em}el\hspace{0.5em}profesor$$

En resumen, la empresa deberá pagar mensualmente 3.454,95 EUR en total. Dado que el proyecto ha durado cuatro meses, el coste total asciende a \textbf{13.819,80 EUR}.

\item \textbf{Hardware}

Los recursos de hardware utilizados han sido un ordenador portátil MacBook Air cuyo coste fue de 1250 EUR y que ha sido amortizado completamente en ejercicios anteriores.
Por lo tanto, podemos concluir que los costes de hardware han sido de \textbf{0 EUR}.

\item \textbf{Software}

En el proyecto, se han utilizado diversas herramientas de software, la mayoría de las cuales son gratuitas. Estas herramientas incluyen:

\begin{itemize}
    \item \textbf{neovim}: Editor de texto gratuito.
    \item \textbf{Navegador web}: Herramienta gratuita para la navegación y pruebas web.
    \item \textbf{CleanShot X}: Herramienta para realizar capturas de pantalla, cuyo coste ha sido de 20 EUR como pago único.
\end{itemize}

Por lo tanto, el coste total del software utilizado en el proyecto ha sido de \textbf{20 EUR}.

\item \textbf{Costes indirectos}

En este apartado se incluyen los costes indirectos asociados al proyecto. Se ha utilizado una tarifa de internet por la que se está pagando 20 EUR mensuales. Dado que el proyecto ha durado cuatro meses, el coste total asciende a:

$$ 20 \text{ EUR/mes} \times 4 \text{ meses} = 80 \text{ EUR} $$

\item \textbf{Total}

La tabla de costes totales se detalla a continuación: \ref{tab:costes-totales}

\begin{table}[ht!]
    \centering
    \resizebox{6.5cm}{!} {
    \begin{tabular}{l c}
    
         \textbf{Tipo de costes}     &  \textbf{Total } \\ \hline
         \textit{Personal}       & 13.819,80 \text{ EUR} \\ 
         \textit{Hardware}    & 0 \text{ EUR} \\  
         \textit{Software}      & 20 \text{ EUR} \\  
         \textit{Costes indirectos} & 80 \text{ EUR} \\ \hline
         \textbf{Total}         & 13.919,80 \text{ EUR}\\ 
    \end{tabular}}
    \caption{Costes totales}
    \label{tab:costes-totales}
\end{table}
\end{itemize}

\subsubsection{Beneficios}
Este es un proyecto con carácter educativo, por lo que no se espera obtener beneficios económicos directos.

\subsection{Viabilidad legal}
A continuación podemos observar en una tabla las licencias que tienen las diferentes dependencias que se encuentran en nuestra aplicación.

\begin{table}[ht!]
    \centering
    \resizebox{15cm}{!} {
    \begin{tabular}{|l|l|l|}
    \hline
         \textbf{Herramienta/Librería} & \textbf{Versión} & \textbf{Licencia} \\ \hline
         {React}       & {17.0.0}  & {MIT} \\ \hline
         {mxGraph}     & {4.1.0}  & {Apache License 2.0} \\ \hline
         {Material UI} & {5.15.19}  & {MIT} \\ \hline
         {Biome}       & {1.7.3}  & {MIT} \\ \hline
         {Playwright}  & {1.44.0}  & {Apache License 2.0} \\ \hline
         {Lefthook}    & {1.6.11}  & {MIT} \\ \hline
         {Vitest}      & {1.6.0}  & {MIT} \\ \hline
    \end{tabular}}
    \caption{Tabla de licencias}
    \label{tab:licencias}
\end{table}

Las herramientas y librerías utilizadas en este proyecto están cubiertas por dos tipos exclusivamente: la licencia MIT\cite{mit:license} y la licencia Apache 2.0\cite{apache:license}. Ambas son permisivas y permiten la reutilización, modificación y distribución del software.

La licencia MIT es muy permisiva y permite a los usuarios utilizar, copiar, modificar, fusionar, publicar, distribuir, sublicenciar y/o vender copias del software \cite{sistemasoperativos:mit}. Siempre que se incluya esta licencia en todas las copias o partes sustanciales del software.

La licencia Apache 2.0 es también muy permisiva a la hora de utilizar, distribuir y modificar versiones del software. También, protege a los desarrolladores de ser responsables de los daños ocasionados \cite{sistemasoperativos:apache}.

Por lo tanto, la licencia de nuestro proyecto puede ser licencia MIT, ya que todas las herramientas y librerías utilizadas permiten su uso bajo estas condiciones. Esto facilita la integración y distribución del proyecto, asegurando que cumpla con las licencias de todas las dependencias utilizadas.

\apendice{Especificación de Requisitos}

\section{Introducción}
En este apartado se enumeran y desarrollan los objetivos y requisitos que debe tener la aplicación según han sido marcados al comienzo del proyecto.

\section{Objetivos generales}
Este proyecto tiene como objetivo los siguientes puntos:
\begin{itemize}
\tightlist
\item Modelar diagramas E-R.
\item Almacenar los diagramas en una estructura interna fácilmente extensible.
\item Validar esa estructura interna para conocer si un diagrama E-R es válido.
\item Exportar un diagrama válido a JSON.
\item Importar y recrear un diagrama en JSON válido.
\item Exportar los diagramas válidos a scripts SQL.
\end{itemize}
\newpage

\section{Catálogo de requisitos}
En este apartado se definirán los requisitos funcionales y no funcionales del proyecto.

\subsection{Requisitos funcionales}
\begin{itemize}
\tightlist
\item \textbf{RF1 - Modelar diagramas E-R:} La aplicación debe permitir a los usuarios crear y manipular diagramas E-R.
\item \textbf{RF2 - Almacenar diagramas:} La aplicación debe almacenar los diagramas en una estructura interna que sea fácilmente extensible.
\item \textbf{RF3 - Validar diagramas:} La aplicación debe validar la estructura interna del diagrama para asegurar su corrección al generar el SQL o exportar/importar JSON.
\item \textbf{RF4 - Exportar diagramas:} La aplicación debe permitir exportar los diagramas validados a un archivo JSON.
\item \textbf{RF5 - Importar diagramas:} La aplicación debe permitir importar diagramas desde un archivo JSON y recrearlos en el entorno de la aplicación.
\item \textbf{RF6 - Generar SQL:} La aplicación debe permitir generar scripts SQL a partir de los diagramas validados.
\end{itemize}

\newpage

\subsection{Requisitos no funcionales}
\begin{itemize}
\tightlist
\item \textbf{RNF1 - Usabilidad:} La aplicación debe ser fácil de usar y tener una interfaz intuitiva para mejorar la experiencia del usuario.
\item \textbf{RNF2 - Rendimiento:} La aplicación debe ser eficiente en términos de rendimiento.
\item \textbf{RNF3 - Compatibilidad:} La aplicación debe ser compatible con diferentes navegadores.
\item \textbf{RNF4 - Seguridad:} La aplicación debe garantizar la seguridad de los datos manejados y almacenados.
\item \textbf{RNF5 - Mantenibilidad:} La aplicación debe ser fácil de mantener y actualizar. Debe poder expandirse en el futuro con nuevos tipos de relaciones y funcionalidades.
\end{itemize}
\clearpage

\section{Especificación de requisitos}

\subsection{Casos de uso}

\begin{table}[p]
	\centering
	\begin{tabularx}{\linewidth}{ p{0.21\columnwidth} p{0.71\columnwidth} }
		\toprule
		\textbf{CU-1}    & \textbf{Modelar diagramas E-R}\\
		\toprule
		\textbf{Versión}              & 1.0    \\
		\textbf{Autor}                & Rubén Maté Iturriaga \\
		\textbf{Requisitos asociados} & RF-1 \\
		\textbf{Descripción}          & Permite al usuario crear y manipular diagramas E-R. \\
		\textbf{Precondición}         & La aplicación debe estar abierta y funcionando. \\
		\textbf{Acciones}             &
		\begin{enumerate}
			\def\labelenumi{\arabic{enumi}.}
			\tightlist
			\item El usuario arrastra y suelta elementos (entidades, relaciones) al canvas.
            \item El usuario añade atributos a las entidades y relaciones (en caso de que fueran N:M).
            \item El usuario edita los nombres y atributos de las entidades y relaciones.
            \item El usuario vincula entidades configurando las relaciones y asignándolas cardinalidades. 
		\end{enumerate}\\
		\textbf{Postcondición}        & El diagrama se almacena en la estructura interna de la aplicación. \\
		\textbf{Excepciones}          & - \\
		\textbf{Importancia}          & Alta \\
		\bottomrule
	\end{tabularx}
	\caption{CU-1 Modelar diagramas E-R.}
\end{table}

\begin{table}[p]
	\centering
	\begin{tabularx}{\linewidth}{ p{0.21\columnwidth} p{0.71\columnwidth} }
		\toprule
		\textbf{CU-2}    & \textbf{Almacenar diagramas}\\
		\toprule
		\textbf{Versión}              & 1.0    \\
	\textbf{Autor}                & Rubén Maté Iturriaga \\
		\textbf{Requisitos asociados} & RF-2 \\
		\textbf{Descripción}          & Permite almacenar los diagramas creados en una estructura interna. \\
		\textbf{Precondición}         & El usuario debe haber creado un diagrama E-R. \\
		\textbf{Acciones}             &
		\begin{enumerate}
			\def\labelenumi{\arabic{enumi}.}
			\tightlist
            \item La estructura interna se mantiene actualizada automáticamente ante los diversos cambios y acciones.
		\end{enumerate}\\
		\textbf{Postcondición}        & El diagrama se almacena correctamente y puede ser recuperado. \\
		\textbf{Excepciones}          & Problemas de almacenamiento (mensaje de error). \\
		\textbf{Importancia}          & Alta \\
		\bottomrule
	\end{tabularx}
	\caption{CU-2 Almacenar diagramas.}
\end{table}

\begin{table}[p]
	\centering
	\begin{tabularx}{\linewidth}{ p{0.21\columnwidth} p{0.71\columnwidth} }
		\toprule
		\textbf{CU-3}    & \textbf{Validar diagramas}\\
		\toprule
		\textbf{Versión}              & 1.0    \\
	\textbf{Autor}                & Rubén Maté Iturriaga \\
		\textbf{Requisitos asociados} & RF-3 \\
		\textbf{Descripción}          & Permite validar la estructura interna del diagrama para asegurar su corrección al generar el SQL o exportar/importar JSON. \\
		\textbf{Precondición}         & El usuario debe haber creado y modelado un diagrama. \\
		\textbf{Acciones}             &
		\begin{enumerate}
			\def\labelenumi{\arabic{enumi}.}
			\tightlist
			\item El usuario selecciona la opción de generar SQL o exportar/importar diagrama.
			\item La aplicación verifica la corrección del diagrama según las reglas E-R.
		\end{enumerate}\\
		\textbf{Postcondición}        & La aplicación muestra un mensaje indicando si el diagrama es válido o no antes de generar el SQL o exportar/importar JSON. \\
		\textbf{Excepciones}          & Mensajes de error con diagnóstico en caso de diagrama no válido. \\
		\textbf{Importancia}          & Alta \\
		\bottomrule
	\end{tabularx}
	\caption{CU-3 Validar diagramas.}
\end{table}

\begin{table}[p]
	\centering
	\begin{tabularx}{\linewidth}{ p{0.21\columnwidth} p{0.71\columnwidth} }
		\toprule
		\textbf{CU-4}    & \textbf{Exportar diagramas}\\
		\toprule
		\textbf{Versión}              & 1.0    \\
	\textbf{Autor}                & Rubén Maté Iturriaga \\
		\textbf{Requisitos asociados} & RF-4 \\
		\textbf{Descripción}          & Permite exportar los diagramas válidos a un archivo JSON. \\
		\textbf{Precondición}         & El diagrama debe ser válido. \\
		\textbf{Acciones}             &
		\begin{enumerate}
			\def\labelenumi{\arabic{enumi}.}
			\tightlist
			\item El usuario selecciona la opción de exportar diagrama.
			\item La aplicación genera un archivo JSON con la estructura interna del diagrama.
			\item El usuario descarga (la descarga se realiza automáticamente al pulsar el botón \texttt{Aceptar}) el archivo JSON.
		\end{enumerate}\\
		\textbf{Postcondición}        & El archivo JSON se guarda en el dispositivo del usuario. \\
		\textbf{Excepciones}          & Problemas de exportación (mensaje de error). \\
		\textbf{Importancia}          & Media \\
		\bottomrule
	\end{tabularx}
	\caption{CU-4 Exportar diagramas.}
\end{table}

\begin{table}[p]
	\centering
	\begin{tabularx}{\linewidth}{ p{0.21\columnwidth} p{0.71\columnwidth} }
		\toprule
		\textbf{CU-5}    & \textbf{Importar diagramas}\\
		\toprule
		\textbf{Versión}              & 1.0    \\
	\textbf{Autor}                & Rubén Maté Iturriaga \\
		\textbf{Requisitos asociados} & RF-5 \\
		\textbf{Descripción}          & Permite importar diagramas desde un archivo JSON y recrearlos en la aplicación. \\
		\textbf{Precondición}         & El usuario debe tener un archivo JSON válido. \\
		\textbf{Acciones}             &
		\begin{enumerate}
			\def\labelenumi{\arabic{enumi}.}
			\tightlist
			\item El usuario selecciona la opción de importar diagrama.
			\item El usuario carga el archivo JSON en la aplicación.
            \item La aplicación hace una validación del diagrama importado y si es correcta recrea el diagrama.
		\end{enumerate}\\
		\textbf{Postcondición}        & El diagrama se visualiza correctamente en la aplicación. \\
		\textbf{Excepciones}          & Problemas de importación (mensaje de error). \\
		\textbf{Importancia}          & Media \\
		\bottomrule
	\end{tabularx}
	\caption{CU-5 Importar diagramas.}
\end{table}

\begin{table}[p]
	\centering
	\begin{tabularx}{\linewidth}{ p{0.21\columnwidth} p{0.71\columnwidth} }
		\toprule
		\textbf{CU-6}    & \textbf{Generar SQL del diagrama modelado}\\
		\toprule
		\textbf{Versión}              & 1.0    \\
		\textbf{Autor}                & Rubén Maté Iturriaga \\
		\textbf{Requisitos asociados} & RF-6 \\
		\textbf{Descripción}          & Permite generar un script SQL a partir del diagrama E-R modelado y validado. \\
		\textbf{Precondición}         & El diagrama debe ser válido. \\
		\textbf{Acciones}             &
		\begin{enumerate}
			\def\labelenumi{\arabic{enumi}.}
			\tightlist
			\item El usuario selecciona la opción de generar SQL.
			\item La aplicación valida el diagrama E-R.
			\item La aplicación genera un script SQL con el paso a tablas del diagrama.
			\item El usuario descarga el archivo SQL (La descarga se realiza automáticamente al pulsar el botón \texttt{Aceptar}).
		\end{enumerate}\\
		\textbf{Postcondición}        & El archivo SQL se guarda en el dispositivo del usuario. \\
		\textbf{Excepciones}          & Problemas de generación de SQL (mensaje de error). \\
		\textbf{Importancia}          & Media \\
		\bottomrule
	\end{tabularx}
	\caption{CU-6 Generar SQL del diagrama modelado.}
\end{table}


\apendice{Especificación de diseño}

\section{Introducción}

En este apéndice se detallan las decisiones de diseño tomadas durante el desarrollo. Se cubren los aspectos relacionados con el diseño de datos, el diseño procedimental y el diseño arquitectónico.

\section{Diseño de datos}
Esta aplicación no tiene como tal una estructura de datos definida puesto que no cuenta con un backend ni base de datos. Sin embargo sí que tiene la estructura JSON interna que contiene nuestro diagrama.

\begin{figure}
\lstset{breaklines=true, basicstyle=\footnotesize}
\begin{lstlisting}[frame=single]
{
    "entities": [
        {
            "idMx": "2",
            "name": "Entidad",
            "position": {
                "x": 254,
                "y": 130
            },
            "attributes": [
                {
                    "idMx": "3",
                    "name": "Atributo",
                    "position": {
                        "x": 374,
                        "y": 130
                    }
                },

            ]
        }
    ],
    "relations": [
    {
      "idMx": "8",
      "name": "Relacion",
      "position": {
        "x": 215,
        "y": 110
      },
      "side1": {
        "idMx": "11",
        "cardinality": "0:1",
        "cell": "11",
        "entity": {
          "idMx": "2"
        }
      },
      "side2": {
        "idMx": "12",
        "cardinality": "1:N",
        "cell": "12",
        "entity": {
          "idMx": "2"
        }
      },
      "canHoldAttributes": false,
      "attributes": []
    }
  ]
}
\end{lstlisting}
\caption{Representación interna del diagrama E-R}
\end{figure}

En esta estructura tenemos 3 entidades definidas que son: Entidades, Relaciones y Atributos.

Su diagrama E-R es el siguiente: \ref{fig:structure-er}
\imagen{structure-er}{Diagrama E-R que representa la estructura interna sobra la que funciona la aplicación.}

\section{Diseño procedimental}
En esta sección se describen los procesos más complejos de la aplicación utilizando diagramas de flujo para entenderlos más fácilmente.

En general ninguno de los procesos de la aplicación requiere grandes diagramas de flujo para comprenderlos. Sin embargo, sí que hay dos funcionalidades que ameritan el tener sus propios diagramas de flujo.

\subsection{Mover objetos}
Es una de las partes más complejas puesto que mover objetos ha de activar diversas tareas según el tipo de objeto que se esté moviendo.
Los objetos pueden ser:
\begin{itemize}
    \item Entidades.
    \item Relaciones.
    \begin{itemize}
        \item Reflexivas.
        \item No reflexivas.
    \end{itemize}
    \item Atributos.
\end{itemize}

\imagen{flow-move-object}{Diagrama de flujo que representa la funcionalidad de mover objetos en el canvas de la aplicación.}

\subsection{Generar el script SQL}
El script SQL es una de las funcionalidades más complejas de la aplicación. Requiere, aparte de la validación, un gran número de pasos y tablas intermedias donde se procesan los diferentes tipos de relaciones.
\imagencontamano{flow-generate-sql}{Diagrama de flujo que representa la funcionalidad de generar el script SQL en la aplicación.}{0.6}


\section{Diseño arquitectónico}
El diseño arquitectónico de la aplicación es bastante simple, puesto que se basa en una arquitectura frontend utilizando React y la librería mxGraph para el modelado de los diagramas E-R.

\subsection{Estructura de paquetes}
La estructura de paquetes de la aplicación se organiza de la siguiente manera:
\begin{itemize}
    \item \textbf{src}: Contiene todo el código fuente de la aplicación.
    \begin{itemize}
        \item \textbf{components/DiagramEditor}
        \begin{itemize}
            \item styles: Carpeta con el css que afecta al componente DiagramEditor
            \item \textbf{DiagramEditor.js}: Componente principal de nuestra aplicación, llama al resto de módulos e inicializa la librería mxGraph.
            \item utils
            \begin{itemize}
                \item \textbf{setInitialConfiguration.js}: Ajusta la configuración la que se inicializa la librería mxGraph.
                \item \textbf{initToolbar.js}: Inicializa la barra de herramientas.
                \item \textbf{addToolbarItem.js}: Función auxiliar que se llama al arrastrar un componente de la barra de tareas al canvas.
                \item configureToolbar.js: Función auxiliar para configurar la barra de herramientas.
                \item getStyleByKey.js: Función auxiliar para obtener estilos de un objeto.
                \item getStyleStringByObj.js: Función auxiliar para obtener estilos de un objeto.
            \end{itemize}
        \end{itemize}
        \item \textbf{utils}:
        \begin{itemize}
            \item \textbf{validation.js}: Módulo que contiene las funciones necesarias para procesar y determina si la estructura interna E-R conforma un diagrama válido.
            \item \textbf{sql.js}: Módulo que contiene las funciones necesarias para procesar y convertir la estructura interna E-R en un script SQL.
        \end{itemize}
    \end{itemize}
\end{itemize}

\subsection{Componentes de la arquitectura}
Aunque la aplicación es puramente frontend, sí se pueden apreciar ciertos componentes externos a la aplicación:
\begin{itemize}
    \item \textbf{Archivos JSON}: Utilizados para la persistencia de los diagramas. Pueden ser importados.
    \item \textbf{LocalStorage}: Almacenamiento local del navegador donde se guarda el estado del diagrama que cuyo modelado esté en curso.
    \item \textbf{Navegador}: Medio a través del cual el usuario interactúa con la aplicación.
\end{itemize}

El siguiente diagrama ilustra estos componentes y su interacción con la aplicación:

\imagen{architecture}{Diagrama con las partes de la aplicación en su ejecución en el navegador.}

\apendice{Documentación técnica de programación}

\section{Introducción}
Este anexo proporciona una guía detallada sobre la estructura de directorios del proyecto, un manual del programador, instrucciones para la compilación, instalación y ejecución del proyecto, así como una descripción de las pruebas del sistema. Está destinado a desarrolladores y posibles contribuidores que deseen comprender mejor el código y colaborar en su desarrollo.

\section{Estructura de directorios}
El proyecto está organizado de la siguiente manera:

\begin{verbatim}
.github/workflows
assets
docs
public
src/
    components/
        DiagramEditor/
            DiagramEditor.js
            utils/
    utils/
        validation.js
        sql.js
tests/
    e2e/
    unit/
        graphs/
\end{verbatim}

\begin{itemize}
    \item \textbf{.github/workflows}: Contiene archivos de configuración para las acciones de github y que establecen los procesos de integración y despliegue continuo.
    \item \textbf{assets}: Archivos de recursos para el proyecto.
    \item \textbf{docs}: Documentación del proyecto donde se encuentran los archivos LateX de esta memoria.
    \item \textbf{public}: Archivos que servirá el servidor web de la aplicación al levantarse.
    \item \textbf{src/components/DiagramEditor}: Contiene la inicialización de mxGraph, el canvas, la toolbar y toda la lógica para modelar el diagrama e interactuar con la aplicación. Hay una explicación un tanto más detallada en el Apéndice C.4.
    \item \textbf{src/utils}: Módulos para validar y exportar el diagrama como SQL.
    \item \textbf{tests/e2e}: Pruebas end-to-end.
    \item \textbf{tests/unit}: Pruebas unitarias.
    \item \textbf{tests/unit/graphs}: Diagramas JSON de ejemplo utilizados en las pruebas unitarias.
\end{itemize}

\section{Manual del programador}
Se proporciona una guía para configurar el proyecto y poder contribuir a su desarrollo.

\subsection{Configuración del entorno de desarrollo}
Para configurar el entorno de desarrollo, se necesitan las siguientes herramientas:
\begin{itemize}
    \item Node.js (versión 18, el repositiorio cuenta con un archivo \texttt{.node-version} que determina la versión node a usar.
    \item npm
    \item git
\end{itemize}

\subsection{Clonación del proyecto}
Uno de los primeros pasos será el de clonar el repositorio remoto a uno local. Preferiblmente se trabajará a través de un fork \cite{github:fork}.
\begin{verbatim}
git clone git@github.com:rubenmate/draw-entity-relation.git
\end{verbatim}

O si se estamos clonando a través del protocolo http (no recomendable, es mejor usar ssh).
\begin{verbatim}
git clone https://github.com/rubenmate/draw-entity-relation.git
\end{verbatim}

\subsection{Instalación de dependencias}
Para instalar las dependencias del proyecto, ejecutar el siguiente comando en la raíz del proyecto:
\begin{verbatim}
npm install
\end{verbatim}

\subsection{Estructura del código}
El código se organiza en módulos para facilitar su mantenimiento y comprensión. Los componentes principales de la aplicación se encuentran en  \textbf{src/components/DiagramEditor}, que contiene toda la lógica para modelar diagramas E-R. Los módulos de validación y generación SQL se encuentran en \textbf{src/utils}.

\subsection{Añadir nuevas funcionalidades}
Para añadir nuevas funcionalidades, hay que seguir estos pasos:
\begin{enumerate}
    \item Crear una rama nueva para su funcionalidad siguiendo la convención de nomenclatura:
    \begin{verbatim}
    git checkout -b <num. del issue>-funcionalidad
    \end{verbatim}
    \item Realizar los cambios necesarios en el código.
    \item Realizar los commits siguiendo la convención \textbf{conventional commits} \cite{conventional-commits}.
    \item Asegurarse de que las pruebas pasen correctamente.
    \item Crear una pull request describiendo los cambios realizados.
\end{enumerate}

Ejemplos de nombres de rama:
\begin{itemize}
    \item Para una nueva funcionalidad: \texttt{29-diagram-persistence}
    \item Para una corrección: \texttt{34-fix-toolbar-bug}
\end{itemize}

Este proyecto cuenta con varias herramientas de integración continua:

Github Actions \cite{github-actions} que se ejecutan automáticamente al hacer commits a la rama principal main o en cualquier rama cuyo objetivo de integración sea la rama principal.

\imagen{github-actions}{Ejecución de las diferentes GitHub Actions en nuestro proyecto.}

Las pruebas pueden ejecutarse también de forma local como se explica en el apartado de \textbf{Pruebas} posterior.

Al realizar un commit se ejecuta un hook que comprueba contra el código incluido en ese commit su corrección semántica según las reglas del linter Biome \cite{biome}.

\section{Compilación, instalación y ejecución del proyecto}
Para compilar, instalar y ejecutar el proyecto, seguir estos pasos:

\subsection{Compilación}
Para compilar el proyecto, ejecutar el siguiente comando:
\begin{verbatim}
npm run build
\end{verbatim}
Con esto se crearía la carpeta build, más información para desplegarla en la sección Despliegue.

\subsection{Ejecución en modo de desarrollo}
Para ejecutar el proyecto en modo de desarrollo, use el siguiente comando:
\begin{verbatim}
npm start
\end{verbatim}
Se levantará el proyecto en la dirección por defecto:
\begin{verbatim}
    http:localhost:3000
\end{verbatim}
Aunque si dicho puerto estaría ocupado se preguntaría al usuario para levantarlo en otro.

\subsection{Despliegue}
Para desplegar el proyecto, se recomienda utilizar Vercel. La configuración es muy básica, tan solo e necesario iniciar sesión en su cuenta de Vercel y conectar con el repositorio del proyecto para que se realice el despliegue de manera automática (inicialmente y con cada nuevo commit).

En caso de querer realizar el despliegue en otra plataforma habría que hacer la build con el comando explicado en el apartado \textbf{Compilación}.
En este punto ya se habría hecho la build y tan solo quedaría servirla con una utilidad como \textbf{serve} \cite{npm:serve}

\begin{enumerate}
    \item Instalar \texttt{serve} a nivel global en el sistema:
    \begin{verbatim}
    npm install -g serve
    \end{verbatim}
    \item Hacer build del proyecto:
    \begin{verbatim}
    npm run build
    \end{verbatim}
    \item Servir la aplicación con \texttt{serve}:
    \begin{verbatim}
    serve -s build
    \end{verbatim}
\end{enumerate}

\section{Pruebas del sistema}
Las pruebas del sistema están divididas en pruebas end-to-end (e2e) y pruebas unitarias.

\subsection{Pruebas unitarias}
Las pruebas unitarias se encuentran en el directorio \textbf{tests/unit}. Estas pruebas verifican la funcionalidad de los componentes individuales del sistema. Los diagramas JSON de ejemplo utilizados en estas pruebas están en \textbf{tests/unit/graphs}.

Para ejecutar las pruebas unitarias, utilizar el siguiente comando:
\begin{verbatim}
npm run test
\end{verbatim}

\imagen{unit-tests}{Ejecución de pruebas unitarias}

\subsection{Pruebas end-to-end}
Las pruebas end-to-end se encuentran en el directorio \textbf{tests/e2e}. Estas pruebas verifican la funcionalidad completa del sistema desde un punto de vista de usuario final.

Para ejecutar las pruebas end-to-end, utilizar el siguiente comando:
\begin{verbatim}
npm run test:e2e
\end{verbatim}

También se pueden lanzar de forma gráfica con \begin{verbatim}
    npx playwright test --ui
\end{verbatim}.

Con este comando se lanza la interfaz gráfica de la herramienta Playwright \cite{playwright} desde donde podemos ver los diferentes tests y ejecutarlos gráficamente.

\imagen{e2e-tests}{Ejecución de pruebas End to End con la interfaz gráfica de Playwright}



\apendice{Documentación de usuario}

\section{Introducción}

\section{Requisitos de usuarios}

\section{Instalación}

\section{Manual del usuario}



\apendice{Anexo de sostenibilización curricular}

\section{Introducción}
Este anexo incluye una reflexión personal sobre los aspectos de sostenibilidad abordados en este Trabajo de Fin de Grado (TFG). Se puede considerar que la sostenibilidad es una idea crucial en la formación académica actual y que además es fundamental para la creación de soluciones tecnológicas que respeten y promuevan el bienestar ambiental, social y económico.

A continuación, se detallan algunas de las competencias de sostenibilidad adquiridas y aplicadas durante el desarrollo del TFG, relacionándolas con los Objetivos de Desarrollo Sostenible (ODS) de la ONU.

\section{Competencias de sostenibilidad adquiridas}

\subsection{ODS 4: Educación de Calidad}
Durante el desarrollo de este TFG, se han adquirido  y aplicado conocimientos avanzados en el ámbito del modelado de diagramas E-R y el desarrollo de aplicaciones web. Un desarrollo implementando buenas prácticas, tales como el uso de metodologías ágiles y herramientas de integración y despliegue continuos, contribuyen a una educación de calidad. Estos nuevos conocimientos no mejora simplemente mis habilidades y conocimientos técnicos, sino que también estimulan un enfoque crítico y reflexivo al abordar la resolución de problemas complejos.

\subsection{ODS 9: Industria, Innovación e Infraestructura}
Este proyecto que se acaba de tratar se centra en la creación de una aplicación web para el modelado de diagramas E-R. Este tipo de herramientas son importantes para la educación, digitalización y modernización de infraestructuras educativas y empresariales. Al facilitar el aprendizaje, diseño y validación de bases de datos; esta aplicación promueve la eficiencia en esta tarea a la vez que se contribuye a la construcción de infraestructuras digitales resilientes y sostenibles.

\subsection{ODS 12: Producción y Consumo Responsables}
El desarrollo de este TFG ha incluido la utilización de recursos de manera eficiente y responsable. Se han usado librerías ya desarrolladas, como son React y mxGraph. Esto reduce la necesidad de recursos adicionales y promueve la reutilización de componentes existentes.

\subsection{ODS 13: Acción por el Clima}
Aunque no se puede decir que este TFG esté directamente relacionado con la mitigación del cambio climático, sí se puede decir que la adopción de prácticas sostenibles en el desarrollo de software tiene un impacto positivo en el medio ambiente.
El uso de herramientas de integración continua y despliegue continuo reduce la necesidad de infraestructuras físicas y, como consecuencia, disminuye la huella de carbono asociada al desarrollo del software.

\subsection{ODS 17: Alianzas para Lograr los Objetivos}
El desarrollo del TFG ha implicado la colaboración con varias partes interesadas, como por ejemplo tutores académicos y otros desarrolladores con problemáticas comunes. Esta colaboración ha fomentado el intercambio de conocimientos, contribuyendo a la construcción de alianzas. El enfoque colaborativo y la comunicación son competencias clave que se han desarrollado y aplicado durante el proyecto, y que también son clave para el desarrollo sostenible.

\section{Reflexión final}
Como conclusión, el desarrollo de este TFG ha permitido la adquisición y aplicación de ciertas competencias de sostenibilidad que están alineadas con varios ODS (Objetivos de Desarrollo Sostenible). 
La implementación de buenas prácticas de desarrollo, la promoción de la innovación y la eficiencia, y el enfoque en la sostenibilidad ambiental y social son aspectos que han tenido presencia en el proyecto de una u otra manera. 
Esta experiencia ha mejorado mis habilidades técnicas. Pero no solo eso, también ha fortalecido mi compromiso con el desarrollo sostenible y responsable en el ámbito tecnológico.

Para más información sobre la sostenibilización curricular, se puede consultar el documento de la CRUE \cite{crue:sostenibilizacion}.




\bibliographystyle{plain}
\bibliography{bibliografiaAnexos}

\end{document}
